% from ega0-1-fr.tex
\oldpage[0\textsubscript{I}]{45}
par cette famille est \emph{surjectif} (autrement dit, si, pour tout
$x\in X$, $\mathcal{F}_x$ est un $\mathcal{O}_x$-module engendré par les
$(s_i)_x$). On dit que $\mathcal{F}$ est \emph{engendré par ses sections
au-dessus de $X$} s'il est engendré par la famille de toutes ces sections (ou
par une sous-famille)~; autrement dit, s'il existe un homomorphisme surjectif
$\mathcal{O}_X^{(I)}\to\mathcal{F}$ pour un $I$ convenable.

On notera qu'un $\mathcal{O}_X$-Module $\mathcal{F}$ peut être tel qu'il
existe un point $x_0\in X$ pour lequel $\mathcal{F}|U$ n'est pas engendré par
ses sections au-dessus de $U$, \emph{quel que soit le voisinage ouvert $U$ de
$x_0$}~: il suffit de prendre $X=\mathbb{R}$, pour $\mathcal{O}_X$ le faisceau
simple $\mathbb{Z}$, pour $\mathcal{F}$ le sous-faisceau algébrique de
$\mathcal{O}_X$ tel que $\mathcal{F}_0=\{0\}$, $\mathcal{F}_x=\mathbb{Z}$ pour
$x\ne0$, et enfin $x_0=0$~: la seule section de $\mathcal{F}|U$ au-dessus de
$U$ est $0$ pour un voisinage $U$ de $0$.
\end{env}

\begin{env}[5.1.2]
\label{0.5.1.2-fr}
Soit $f:X\to Y$ un morphisme d'espaces annelés. Si $\mathcal{F}$ est un
$\mathcal{O}_X$-Module engendré par ses sections au-dessus de $X$, alors
l'homomorphisme canonique
$f^*(f_*(\mathcal{F}))\to\mathcal{F}$ (4.4.3.3) est \emph{surjectif}~: en
effet, avec les notations de (5.1.1), $s_i\otimes1$ est une section de
$f^*(f_*(\mathcal{F}))$ au-dessus de $X$, et son image dans $\mathcal{F}$ est
$s_i$. L'exemple de (5.1.1) où $f$ est l'identité, montre que la réciproque de
cette proposition est inexacte en général.
\end{env}

\begin{env}[5.1.3]
\label{0.5.1.3-fr}
On dit qu'un $\mathcal{O}_X$-Module $\mathcal{F}$ est \emph{quasi-cohérent}
si, pour tout $x\in X$, il y a un voisinage ouvert $U$ de $x$ tel que
$\mathcal{F}|V$ soit isomorphe au \emph{conoyau} d'un homomorphisme de la
forme
$\mathcal{O}_X^{(I)}|V\to\mathcal{O}_X^{(J)}|V$, où $I$ et $J$ soient des
ensembles d'indices arbitraires. Il est clair que $\mathcal{O}_X$ lui-même est
un $\mathcal{O}_X$-Module quasi-cohérent, et que toute somme directe de
$\mathcal{O}_X$-Modules quasi-cohérents est un $\mathcal{O}_X$-Module
quasi-cohérent. On dit qu'une $\mathcal{O}_X$-Algèbre $\mathcal{A}$ est
\emph{quasi-cohérente} si elle l'est en tant que $\mathcal{O}_X$-Module.
\end{env}

\begin{env}[5.1.4]
\label{0.5.1.4-fr}
Soit $f:X\to Y$ un morphisme d'espaces annelés. Si $\mathcal{G}$ est un
$\mathcal{O}_Y$-Module quasi-cohérent, alors $f^*(\mathcal{G})$ est un
$\mathcal{O}_X$-Module quasi-cohérent. En effet, pour tout $x\in X$, il y a
un voisinage ouvert $V$ de $f(x)$ dans $Y$ tel que $\mathcal{G}|V$ soit le
conoyau d'un homomorphisme
$\mathcal{O}_Y^{(I)}|V\to\mathcal{O}_Y^{(J)}|V$. Si $U=f^{-1}(V)$, et si
$f_U$ est la restriction de $f$ à $U$, on a
$f^*(\mathcal{G})|U=f_U^*(\mathcal{G}|V)$~; comme $f_U^*$ est exact à droite
et commute aux sommes directes, $f_U^*(\mathcal{G}|V)$ est conoyau d'un
homomorphisme
$\mathcal{O}_X^{(I)}|U\to\mathcal{O}_X^{(J)}|U$.
\end{env}

\subsection{Faisceaux de type fini}
\label{subsection:0.5.2-fr}

\begin{env}[5.2.1]
\label{0.5.2.1-fr}
On dit qu'un $\mathcal{O}_X$-Module $\mathcal{F}$ est \emph{de type fini} si,
pour tout $x\in X$, il existe un voisinage ouvert $U$ de $x$ tel que
$\mathcal{F}|U$ soit engendré par une famille \emph{finie} de sections
au-dessus de $U$, ou encore soit isomorphe à un faisceau quotient d'un
faisceau de la forme $(\mathcal{O}_X|U)^p$ où $p$ est fini. Tout faisceau
quotient d'un faisceau de type fini est de type fini, ainsi que toute somme
directe finie et tout produit tensoriel fini de faisceaux de type fini. Un
$\mathcal{O}_X$-Module de type fini n'est pas nécessairement quasi-cohérent,
comme on peut le voir pour le $\mathcal{O}_X$-Module
$\mathcal{O}_X/\mathcal{F}$, où $\mathcal{F}$ est l'exemple de (5.1.1). Si
$\mathcal{F}$ est de type fini, $\mathcal{F}_x$ est un
$\mathcal{O}_x$-module de type fini pour tout $x\in X$, mais l'exemple
(5.1.1) montre que cette condition nécessaire n'est pas en général
suffisante.
\end{env}

\begin{env}[5.2.2]
\label{0.5.2.2-fr}
Soit $\mathcal{F}$ un $\mathcal{O}_X$-Module \emph{de type fini}. Si $s_i$
$(1\leq i\leq n)$ sont des sections de $\mathcal{F}$ au-dessus d'un voisinage
ouvert $U$ d'un point $x\in X$ et si les $(s_i)_x$ engendrent
$\mathcal{F}_x$, il existe un voisinage ouvert $V\subset U$ de $x$ tel que
les $(s_i)_y$ engendrent $\mathcal{F}_y$ pour tout $y\in V$ (FAC, I, 2, 12,
prop.~1). On en conclut en particulier que le support de $\mathcal{F}$ est
\emph{fermé}.

